% together/together.tex
% mainfile: ../perfbook.tex
% SPDX-License-Identifier: CC-BY-SA-3.0

\QuickQuizChapter{chp:Putting It All Together}{融会贯通}{qqztogether}
%
\Epigraph{你不是先学射击，然后学发射，再学控制旋转——你学的是
	  发射-射击-旋转。}{\emph{Ender's Shadow}，Orson Scott Card}

% And the paragraph preceding this is also instructive:
% ``I may be pissed off, but that doesn't mean I can't learn.''

本\lcnamecref{chp:Putting It All Together}
给出了一些关于并发编程难题的提示。
\Cref{sec:together:Counter Conundrums}
考虑计数器难题，
\cref{sec:together:Refurbish Reference Counting}
翻新引用计数，
\cref{sec:together:Hazard-Pointer Helpers}
帮助解决 hazard pointer 问题，
\cref{sec:together:Sequence-Locking Specials}
探讨 sequence lock 的特殊用法，
\cref{sec:together:RCU Rescues}
则反思 RCU 的救援之道。
最后，
尽管最佳的性能和可扩展性来自于设计而非事后的微优化，
但微优化对于追求绝对最佳的性能和可扩展性仍然是必要的。
因此，
\cref{sec:together:Micro-Optimization}
将漫谈一些微优化技巧。

% together/count.tex
% mainfile: ../perfbook.tex
% SPDX-License-Identifier: CC-BY-SA-3.0

\section{计数器难题}
\label{sec:together:Counter Conundrums}
%
\epigraph{Ford继续安静地计数。
	  这大概是你能对计算机做的最具攻击性的事情了，
	  相当于走到一个人面前说
	  ``血……血……血……血……''}
	 {Douglas Adams}

本\lcnamecref{sec:together:Counter Conundrums}
概述计数器难题的解决方案。

\subsection{计算更新次数}
\label{sec:together:Counting Updates}

假设 Schr\"odinger（参见
\cref{sec:datastruct:Motivating Application}）
想要计算每个动物的更新次数，
并且这些更新使用每数据元素锁来同步。
如何最好地进行这种计数？

当然，\cref{chp:Counting} 中的任何计数算法都可能适用，
但最优的方法相当简单。
只需在每个数据元素中放置一个计数器，
并在该元素锁的保护下递增它！

如果读者无锁（lock-free）地访问计数，
那么更新者应使用 \co{WRITE_ONCE()} 来更新计数器，
而无锁读者应使用 \co{READ_ONCE()} 来加载它。

\subsection{计算查找次数}
\label{sec:together:Counting Lookups}

假设 Schr\"odinger 还想计算每个动物的查找次数，
其中查找由 RCU 保护。
如何最好地进行这种计数？

一种方法是使用每元素锁来保护查找计数器，
如 \cref{sec:together:Counting Updates} 中所述。
不幸的是，这将要求所有查找都获取此锁，
这在大型系统上会成为严重的瓶颈。

另一种方法是"直接拒绝"计数，
效仿 \co{noatime} 挂载选项的做法。
如果这种方法可行，它显然是最好的：
毕竟，没有什么比什么都不做更快。
如果查找计数不可或缺，请继续阅读！

\cref{chp:Counting} 中的任何计数器都可以派上用场，
其中 \cref{sec:count:Statistical Counters} 描述的统计计数器（statistical counter）
可能是最常见的选择。
然而，这会导致较大的内存占用：
所需的计数器数量是数据元素数量乘以线程（thread）数量。

如果这种内存开销过大，
一种方法是保持每核心甚至每插槽的计数器而非每 CPU 计数器，
同时参考
\cref{fig:datastruct:Read-Only Hash-Table Performance For Schroedinger's Zoo; 448 CPUs}
中所示的 hash table 性能结果。
这将要求计数器递增是原子操作（atomic operation），
特别是对于用户态执行，
因为给定的线程可能在任何时候迁移到另一个 CPU。

如果某些元素被非常频繁地查找，
有许多方法可以通过维护每线程日志来批量更新，
其中同一元素的多个日志条目可以合并。
在给定的日志条目具有足够大的增量或经过足够的时间后，
日志条目可以应用到相应的数据元素上。
Silas Boyd-Wickizer 在形式化这一概念方面做了一些工作~\cite{SilasBoydWickizerPhD}。

% together/refcnt.tex
% mainfile: ../perfbook.tex
% SPDX-License-Identifier: CC-BY-SA-3.0

\section{翻新引用计数}
\label{sec:together:Refurbish Reference Counting}
%
\epigraph{计数是这一代人的信仰。
	  它是他们的希望和救赎。}
	 {Gertrude Stein}

虽然引用计数在概念上是一种简单的技术，
但将其应用于并发软件时，许多魔鬼隐藏在细节之中。
毕竟，如果对象不会被提前释放，
就根本不需要
\IXalt{reference counter}{reference count}。
但是，如果对象可以被释放，
那么在引用获取过程本身中，
又是什么阻止了释放的发生呢？

有多种方法可以翻新引用计数以用于并发软件，包括：

\begin{enumerate}
\item	在操作引用计数时，必须持有位于对象外部的锁。
\item	对象创建时具有非零引用计数，
	只有当引用计数器的当前值非零时才可获取新引用。
	如果某线程没有对给定对象的引用，
	它可以向已持有引用的另一个线程寻求帮助。
\item	在某些情况下，hazard pointer 可以作为引用计数器的
	直接替代品使用。
\item	为对象提供\IX{existence guarantee}，
	从而防止在其他实体可能正在尝试获取引用时被释放。
	存在性保证通常由自动垃圾收集器提供，
	而且如
	\cref{sec:defer:Hazard Pointers,sec:defer:Read-Copy Update (RCU)}
	所示，分别由 hazard pointer 和 RCU 提供。
\item	为对象提供\IXalt{type-safety guarantee}{type-safe memory}。
	一旦获取了引用，还必须执行额外的身份检查。
	类型安全保证可以由专用内存分配器提供，
	例如 Linux 内核中的
	\co{SLAB_TYPESAFE_BY_RCU} 特性，
	如 \cref{sec:defer:Read-Copy Update (RCU)} 所示。
\end{enumerate}

当然，任何提供存在性保证的机制，
根据定义也提供类型安全保证。
这产生了四大类引用获取保护方式：
引用计数、hazard pointer、sequence lock 和 RCU\@。

\QuickQuiz{
	为什么不使用简单的 \IXacrml{cas} 操作来实现引用获取，
	使其仅在引用计数器非零时才获取引用？
}\QuickQuizAnswer{
	虽然这可以解决最后一个引用的释放与新引用的获取之间的竞争，
	但它完全无法防止数据结构被释放并重新分配，
	可能作为完全不同类型的结构。
	如果将``简单的 \acrml{cas} 操作''应用于
	不同类型的结构，很可能会产生未定义的结果。

	简而言之，使用 \acrml{cas} 等原子操作
	绝对需要类型安全保证或存在性保证。

	但如果确实有必要让类型发生变化呢？

	一种方法是让每种此类类型将引用计数器放在相同的位置，
	这样只要重新分配的结果是该类型组中的对象，一切就没问题。
	如果你用 C 语言这样做，
	请确保在每个出现引用计数器的结构中添加注释。
	在 C++ 中，使用继承和模板。
}\QuickQuizEnd

\begin{table}
\renewcommand*{\arraystretch}{1.25}
\rowcolors{3}{}{lightgray}
\small
\centering
\begin{tabular}{lcccc}
	\toprule
	& \multicolumn{4}{c}{释放} \\
	\cmidrule(l){2-5}
	获取	& 锁
				& \parbox[c]{.5in}{引用\\计数}
					& \parbox[c]{.5in}{Hazard\\Pointer}
						& RCU \\
	\cmidrule{1-1} \cmidrule(l){2-5}
	锁		& $-$	& CAM	& M	& CA  \\
	\parbox[c][6ex]{.6in}{引用\\计数}
			& A	& AM    & M	& A   \\
	\parbox[c][6ex]{.6in}{Hazard\\Pointer}
			& M	& M	& M	& M   \\
	RCU		& CA	& MCA	& M	& CA  \\
	\bottomrule
\end{tabular}
\caption{引用计数的同步方式}
\label{tab:together:Synchronizing Reference Counting}
\end{table}

鉴于引用计数的关键问题是
引用获取与对象释放之间的同步，
我们有九种可能的机制组合，如
\cref{tab:together:Synchronizing Reference Counting} 所示。
该表将引用计数机制分为以下几大类：
\begin{enumerate}
\item	简单计数，既不需要原子操作，
	也不需要\IXpl{memory barrier}，也不需要对齐约束（``$-$''）。
\item	无内存屏障的原子计数（``A''）。
\item	原子计数，仅在释放时需要内存屏障（``AM''）。
\item	原子计数，获取操作中结合了检查的原子操作，
	仅在释放时需要内存屏障（``CAM''）。
\item	原子计数，获取操作中结合了检查的原子操作（``CA''）。
\item	简单计数，结合了检查和完整内存屏障（``M''）。
\item	原子计数，获取操作中结合了检查的原子操作，
	且在获取时也需要内存屏障（``MCA''）。
\end{enumerate}
然而，由于所有返回值的 Linux 内核原子操作
都被定义为包含内存屏障，\footnote{
	\co{atomic_read()} 和 \co{ATOMIC_INIT()} 是
	证明规则的例外。}
所有释放操作都包含内存屏障，
所有带检查的获取操作也包含内存屏障。
因此，``CA'' 和 ``MCA'' 等价于 ``CAM''，
所以下面只有前四种情况和第六种情况的小节：
``$-$''、``A''、``AM''、``CAM'' 和 ``M\@''。
后续小节描述了当引用获取和释放非常频繁，
而引用计数仅在极少数情况下需要检查是否为零时，
可以提高性能的优化方法。

\subsection{引用计数类别的实现}
\label{sec:together:Implementation of Reference-Counting Categories}

由锁保护的简单计数（``$-$''）在
\cref{sec:together:Simple Counting} 中描述，
无内存屏障的原子计数（``A''）在
\cref{sec:together:Atomic Counting} 中描述，
带获取内存屏障的原子计数（``AM''）在
\cref{sec:together:Atomic Counting With Release Memory Barrier} 中描述，
而
带检查和释放内存屏障的原子计数（``CAM''）在
\cref{sec:together:Atomic Counting With Check and Release Memory Barrier} 中描述。
Hazard pointer 的使用在
\cref{sec:defer:Hazard Pointers}
（\cpageref{sec:defer:Hazard Pointers}）
以及
\cref{sec:together:Hazard-Pointer Helpers} 中描述。

\subsubsection{简单计数}
\label{sec:together:Simple Counting}

简单计数，既不需要原子操作也不需要内存屏障，
可以在引用计数器的获取和释放
都由同一个锁保护时使用。
在这种情况下，显然引用计数本身
可以非原子地操作，因为锁提供了所有
必要的互斥、内存屏障、原子指令和
编译器优化禁用。
当锁不仅需要保护引用计数，还需要保护其他操作，
但在锁释放后必须持有对象的引用时，
这是首选方法。
\Cref{lst:together:Simple Reference-Count API} 展示了
一个可能用于实现简单非原子引用计数的简单 API——
尽管简单引用计数几乎总是以内联方式编码。

\begin{listing}
\begin{fcvlabel}[ln:together:Simple Reference-Count API]
\begin{VerbatimL}[commandchars=\\\[\]]
struct sref {
	int refcount;
};

void sref_init(struct sref *sref)
{
	sref->refcount = 1;
}

void sref_get(struct sref *sref)
{
	sref->refcount++;
}

int sref_put(struct sref *sref,
             void (*release)(struct sref *sref))
{
	WARN_ON(release == NULL);
	WARN_ON(release == (void (*)(struct sref *))kfree);

	if (--sref->refcount == 0) {
		release(sref);
		return 1;
	}
	return 0;
}
\end{VerbatimL}
\end{fcvlabel}
\caption{简单引用计数 API}
\label{lst:together:Simple Reference-Count API}
\end{listing}

\subsubsection{原子计数}
\label{sec:together:Atomic Counting}

简单原子计数可用于任何获取引用的 CPU
必须已经持有引用的情况。
这种方式用于单个 CPU 为其私有用途创建对象，
但必须允许来自其他 CPU、任务、定时器处理程序等的访问。
任何将对象传递出去的 CPU 必须首先代表接收方获取一个新引用，
或者在传递后停止进一步的访问。
在 Linux 内核中，\co{kref} 原语用于实现
这种引用计数方式，如
\cref{lst:together:Linux Kernel kref API} 所示。\footnote{
	截至 Linux v4.10。
	Linux v4.11 引入了 \co{refcount_t} API，
	提高了弱序平台上的效率，
	但在功能上等同于它所替代的 \co{atomic_t}。}

在这种情况下需要原子计数，
因为锁不保护所有引用计数操作，
这意味着两个不同的 CPU 可能同时操作引用计数。
如果使用普通的递增和递减，一对 CPU 可能同时获取引用计数，
也许都获得值 ``3''。
如果两者都递增其值，它们都将获得 ``4''，
并且都将此值存储回计数器。
由于计数器的新值应该是 ``5''，
其中一次递增就丢失了。
因此，计数器递增和递减都必须使用原子操作。

如果释放由锁、hazard pointer 或 RCU 保护，
则\emph{不}需要内存屏障，但原因各不相同。
在锁的情况下，锁提供所需的内存屏障
（以及编译器优化禁用），
锁还防止一对释放操作并发运行。
在 hazard pointer 和 RCU 的情况下，清理将被延迟，
所需的内存屏障或编译器优化禁用
将由 hazard pointer 或 RCU 基础设施提供。
因此，如果两个 CPU 并发释放最后两个引用，
实际的清理将被推迟到两个 CPU 都释放了 hazard pointer
或分别退出了 RCU 读侧临界区之后。

\QuickQuiz{
	为什么不需要防范一个 CPU 在另一个 CPU
	释放最后一个引用后立即获取引用的情况？
}\QuickQuizAnswer{
	因为 CPU 必须已经持有引用才能合法地获取另一个引用。
	因此，如果一个 CPU 释放了最后一个引用，
	最好不要有任何 CPU 正在获取新引用！
}\QuickQuizEnd

\begin{listing}
\begin{fcvlabel}[ln:together:Linux Kernel kref API]
\begin{VerbatimL}[commandchars=\\\[\]]
struct kref {						\lnlbl[kref:b]
	atomic_t refcount;
};							\lnlbl[kref:e]

void kref_init(struct kref *kref)			\lnlbl[init:b]
{
	atomic_set(&kref->refcount, 1);
}							\lnlbl[init:e]

void kref_get(struct kref *kref)			\lnlbl[get:b]
{
	WARN_ON(!atomic_read(&kref->refcount));
	atomic_inc(&kref->refcount);
}							\lnlbl[get:e]

static inline int					\lnlbl[sub:b]
kref_sub(struct kref *kref, unsigned int count,
         void (*release)(struct kref *kref))
{
	WARN_ON(release == NULL);

	if (atomic_sub_and_test((int) count,		\lnlbl[check]
	                        &kref->refcount)) {
		release(kref);				\lnlbl[rel]
		return 1;				\lnlbl[ret:1]
	}
	return 0;
}							\lnlbl[sub:e]
\end{VerbatimL}
\end{fcvlabel}
\caption{Linux 内核 \tco{kref} API}
\label{lst:together:Linux Kernel kref API}
\end{listing}

\begin{fcvref}[ln:together:Linux Kernel kref API]
\co{kref} 结构本身由单个原子数据项组成，
如 \cref{lst:together:Linux Kernel kref API}
的 \clnrefrange{kref:b}{kref:e} 所示。
\co{kref_init()} 函数在 \clnrefrange{init:b}{init:e}
将计数器初始化为值 ``1''。
注意 \co{atomic_set()} 原语是一个简单的赋值，
其名称源于 \co{atomic_t} 的数据类型
而非操作本身。
\co{kref_init()} 函数必须在对象创建期间调用，
即在对象对任何其他 CPU 可用之前调用。

\co{kref_get()} 函数在 \clnrefrange{get:b}{get:e}
无条件地原子递增计数器。
\co{atomic_inc()} 原语不一定在所有平台上
显式禁用编译器优化，但 \co{kref}
原语在独立模块中且 Linux 内核构建过程不进行跨模块优化，
所以效果相同。

\co{kref_sub()} 函数在 \clnrefrange{sub:b}{sub:e}
原子递减计数器，如果结果为零，
\clnref{rel} 调用指定的 \co{release()} 函数，
\clnref{ret:1} 返回，通知调用者 \co{release()} 已被调用。
否则，\co{kref_sub()} 返回零，
通知调用者 \co{release()} 未被调用。
\end{fcvref}

\QuickQuizSeries{%
\QuickQuizB{
	\begin{fcvref}[ln:together:Linux Kernel kref API]
	假设在 \cref{lst:together:Linux Kernel kref API}
	的 \clnref{check} 上调用 \co{atomic_sub_and_test()} 之后，
	某个其他 CPU 调用了 \co{kref_get()}。
	这不就导致那个 CPU 持有了
	对已释放对象的非法引用吗？
	\end{fcvref}
}\QuickQuizAnswerB{
	如果这些函数被正确使用，这是不可能发生的。
	除非你已经持有引用，否则调用 \co{kref_get()} 是非法的，
	在这种情况下，\co{kref_sub()} 不可能将计数器递减到零。
}\QuickQuizEndB
%
\QuickQuizM{
	假设 \co{kref_sub()} 返回零，
	表示 \co{release()} 函数未被调用。
	在什么条件下调用者可以依赖包含对象的继续存在？
}\QuickQuizAnswerM{
	除非调用者知道至少会有一个引用继续存在，
	否则不能依赖对象的继续存在。
	通常，调用者无法知道这一点，
	因此必须在调用 \co{kref_sub()} 后
	小心避免引用该对象。

	鼓励有兴趣的读者使用 RCU 来解决这一限制，
	特别是使用 \co{call_rcu()}。
}\QuickQuizEndM
%
\QuickQuizE{
	为什么不直接将 \co{kfree()} 作为释放函数传入？
}\QuickQuizAnswerE{
	因为 \co{kref} 结构通常嵌入在更大的结构中，
	需要释放整个结构，而不仅仅是 \co{kref} 字段。
	这通常通过定义一个包装函数来完成，
	该函数执行 \co{container_of()} 然后执行 \co{kfree()}。
}\QuickQuizEndE
}

\subsubsection{带释放内存屏障的原子计数}
\label{sec:together:Atomic Counting With Release Memory Barrier}

带释放内存屏障的原子引用计数被 Linux 内核的网络层
用于跟踪数据包路由中使用的目标缓存。
实际实现要复杂得多；本节重点关注
\co{struct dst_entry} 引用计数处理中
与此用例匹配的方面，
如 \cref{lst:together:Linux Kernel dst-clone API} 所示。\footnote{
	截至 Linux v4.13。
	Linux v4.14 增加了一层间接引用以允许
	更全面的调试检查，但在没有 bug 的情况下
	整体效果是相同的。}

\begin{listing}
\begin{fcvlabel}[ln:together:Linux Kernel dst-clone API]
\begin{VerbatimL}[commandchars=\\\[\]]
static inline
struct dst_entry * dst_clone(struct dst_entry * dst)
{
	if (dst)
		atomic_inc(&dst->__refcnt);
	return dst;
}

static inline
void dst_release(struct dst_entry * dst)
{
	if (dst) {
		WARN_ON(atomic_read(&dst->__refcnt) < 1);
		smp_mb__before_atomic_dec();		\lnlbl[mb]
		atomic_dec(&dst->__refcnt);
	}
}
\end{VerbatimL}
\end{fcvlabel}
\caption{Linux 内核 \tco{dst_clone} API}
\label{lst:together:Linux Kernel dst-clone API}
\end{listing}

\co{dst_clone()} 原语可在调用者已持有对指定 \co{dst_entry}
的引用时使用，此时它获取另一个引用，
该引用可以传递给内核中的其他实体。
因为调用者已持有引用，
\co{dst_clone()} 不需要执行任何内存屏障。
将 \co{dst_entry} 传递给其他实体可能
需要也可能不需要内存屏障，但如果需要，
它将嵌入在用于传递 \co{dst_entry} 的机制中。

\begin{fcvref}[ln:together:Linux Kernel dst-clone API]
\co{dst_release()} 原语可从任何环境调用，
调用者很可能在调用 \co{dst_release()} 之前
立即引用 \co{dst_entry} 结构的元素。
因此 \co{dst_release()} 原语在 \clnref{mb}
包含内存屏障，防止编译器和 CPU 错误排序访问。
\end{fcvref}

请注意，使用 \co{dst_clone()} 和
\co{dst_release()} 的程序员不需要了解内存屏障，
只需了解使用这两个原语的规则。

\subsubsection{带检查和释放内存屏障的原子计数}
\label{sec:together:Atomic Counting With Check and Release Memory Barrier}

考虑这样一种情况：调用者必须能够获取对
其尚未持有引用的对象的新引用，
但该对象的存在是有保证的。
初始引用计数获取现在可以与引用计数释放并发运行，
这增加了进一步的复杂性。
假设引用计数释放发现引用计数的新值为零，
表明现在可以安全地清理引用计数对象。
显然，我们不能允许在清理开始后再进行引用计数获取，
因此获取必须包含对零引用计数的检查。
此检查必须是原子递增操作的一部分，
如下所示。

\QuickQuiz{
	为什么不能在简单的 \qco{if} 语句中
	检查零引用计数，并在其 \qco{then} 子句中
	进行原子递增？
}\QuickQuizAnswer{
	假设 \qco{if} 条件完成，
	发现引用计数器的值等于 1。
	假设一个释放操作执行，
	将引用计数器递减到零，因此开始清理操作。
	但现在 \qco{then} 子句可以将计数器递增回值 1，
	允许在对象被清理后继续使用。

	这个清理后使用的 bug 与
	完整的释放后使用 bug 一样严重。
}\QuickQuizEnd

Linux 内核的 \co{fget()} 和 \co{fput()} 原语
使用这种引用计数方式。
这些函数的简化版本如
\cref{lst:together:Linux Kernel fget/fput API} 所示。\footnote{
	截至 Linux v2.6.38。
	v2.6.39 添加了额外的 \co{O_PATH} 功能，
	v3.14 进行了重构，v4.1 减少了 \co{mmap_sem} 竞争。}

\begin{listing}
\begin{fcvlabel}[ln:together:Linux Kernel fget/fput API]
\begin{VerbatimL}[commandchars=\\\@\$]
struct file *fget(unsigned int fd)
{
	struct file *file;
	struct files_struct *files = current->files;	\lnlbl@fetch$


	rcu_read_lock();				\lnlbl@rrl$
	file = fcheck_files(files, fd);			\lnlbl@lookup$
	if (file) {
		if (!atomic_inc_not_zero(&file->f_count)) { \lnlbl@acq$
			rcu_read_unlock();		\lnlbl@rru1$
			return NULL;			\lnlbl@ret:null$
		}
	}
	rcu_read_unlock();				\lnlbl@rru2$
	return file;					\lnlbl@ret:file$
}

struct file *
fcheck_files(struct files_struct *files, unsigned int fd)
{
	struct file * file = NULL;
	struct fdtable *fdt = rcu_dereference((files)->fdt);  \lnlbl@deref:fdt$

	if (fd < fdt->max_fds)				\lnlbl@range$
		file = rcu_dereference(fdt->fd[fd]);	\lnlbl@deref:fd$
	return file;					\lnlbl@ret:file2$
}

void fput(struct file *file)
{
	if (atomic_dec_and_test(&file->f_count))	\lnlbl@dec$
		call_rcu(&file->f_u.fu_rcuhead, file_free_rcu);  \lnlbl@call$
}

static void file_free_rcu(struct rcu_head *head)
{
	struct file *f;

	f = container_of(head, struct file, f_u.fu_rcuhead); \lnlbl@obtain$
	kmem_cache_free(filp_cachep, f);		\lnlbl@free$
}
\end{VerbatimL}
\end{fcvlabel}
\caption{Linux 内核 \tco{fget}/\tco{fput} API}
\label{lst:together:Linux Kernel fget/fput API}
\end{listing}

\begin{fcvref}[ln:together:Linux Kernel fget/fput API]
\co{fget()} 的 \clnref{fetch} 获取指向当前进程文件描述符表的指针，
该表很可能与其他进程共享。
\Clnref{rrl} 调用 \co{rcu_read_lock()}，
进入 RCU 读侧临界区。
任何后续 \co{call_rcu()} 原语的回调函数
将被推迟到匹配的 \co{rcu_read_unlock()} 到达时
（在此示例中为 \clnref{rru1} 或 \lnref{rru2}）。
\Clnref{lookup} 查找与 \co{fd} 参数指定的文件描述符
对应的文件结构，稍后将进行描述。
如果存在与指定文件描述符对应的打开文件，
则 \clnref{acq} 尝试原子地获取引用计数。
如果失败，\clnrefrange{rru1}{ret:null} 退出 RCU 读侧临界区
并报告失败。
否则，如果尝试成功，\clnrefrange{rru2}{ret:file}
退出读侧临界区并返回指向文件结构的指针。

\co{fcheck_files()} 原语是 \co{fget()} 的辅助函数。
\Clnref{deref:fdt} 使用 \co{rcu_dereference()} 安全地获取
指向此任务当前文件描述符表的 RCU 保护指针，
\clnref{range} 检查指定的文件描述符是否在范围内。
如果是，\clnref{deref:fd} 获取指向文件结构的指针，
同样使用 \co{rcu_dereference()} 原语。
\Clnref{ret:file2} 然后返回指向文件结构的指针，
或者在失败时返回 \co{NULL}。

\co{fput()} 原语释放对文件结构的引用。
\Clnref{dec} 原子递减引用计数，如果结果为零，
\clnref{call} 调用 \co{call_rcu()} 原语
以释放文件结构（通过 \co{call_rcu()} 第二个参数中
指定的 \co{file_free_rcu()} 函数），但仅在
所有当前正在执行的 RCU 读侧临界区完成之后，
即在一个 RCU \IX{grace period} 过去之后。

一旦宽限期完成，\co{file_free_rcu()} 函数
在 \clnref{obtain} 获取指向文件结构的指针，
并在 \clnref{free} 释放它。
\end{fcvref}

这段代码片段展示了如何使用 RCU 来保证
对象内引用计数递增期间的存在性。

\subsection{计数器优化}
\label{sec:together:Counter Optimizations}

在某些情况下，递增和递减很常见，
但检查零的情况很少，
维护每 CPU 或每任务计数器是有意义的，
如 \cref{chp:Counting} 中所讨论的。
例如，参见关于可睡眠读-复制-更新（SRCU）的论文，
该论文将此技术应用于 RCU~\cite{PaulEMcKenney2006c}。
这种方法消除了递增和递减原语中
对原子指令或内存屏障的需要，
但仍然要求禁用代码移动编译器优化。
此外，检查聚合引用计数是否达到零的原语
（如 \co{synchronize_srcu()}）可能相当慢。
这凸显了这些技术是为引用频繁获取和释放，
但很少需要检查零引用计数的情况而设计的。

然而，通常情况下使用引用计数需要
对原本只读的数据结构进行写入（通常是原子写入）。
在这种情况下，引用计数给读者带来了
昂贵的缓存未命中。

因此，值得研究不需要读者
写入正在遍历的数据结构的同步机制。
一种可能是 \cref{sec:defer:Hazard Pointers}
中介绍的 hazard pointer，
另一种是 \cref{sec:defer:Read-Copy Update (RCU)}
中介绍的 RCU。

% together/hazptr.tex
% mainfile: ../perfbook.tex
% SPDX-License-Identifier: CC-BY-SA-3.0

\section{风险指针助手}
\label{sec:together:Hazard-Pointer Helpers}
%
\epigraph{重要的是小事情，成百上千的小事情。}
	 {Cliff Shaw}

本节探讨一些可以借助 hazard pointer 解决的问题。
此外，hazard pointer 有时也可以用来解决
\cref{sec:together:RCU Rescues} 中提出的问题，反之亦然。

\subsection{可扩展引用计数}
\label{sec:together:Scalable Reference Count}

假设引用计数正在成为性能或可扩展性的瓶颈。
你能做什么？

一种方法是改用\IXpl{hazard pointer}。

两者之间存在一些差异，
最显著的也许是使用 hazard pointer 时，
确定相应的引用计数何时达到零的代价极其昂贵。

解决此问题的一种方法是在引用计数和 hazard pointer 之间分担负载。
每个数据元素都有一个引用计数器，
一方面跟踪引用该元素的其他数据元素的数量，
另一方面读者使用 hazard pointer。

要使这种安排既高效又正确可能相当具有挑战性，
因此有兴趣的读者可以查看
Folly 开源库中实现的 UnboundedQueue 和 ConcurrentHashMap 数据结构。\footnote{
	\url{https://github.com/facebook/folly}}

\subsection{长时间访问}
\label{sec:together:Long-Duration Accesses}

假设一个读写锁的读者持有锁的时间太长，
导致更新被过度延迟。
如果该读者可以合理地转换为使用引用计数
而非读写锁，但性能和可扩展性的考虑又阻止使用
实际的引用计数器，那么 hazard pointer 提供了
引用计数的一种可扩展变体。

关键点在于，读写锁的读者会阻塞该锁的所有更新，
而 hazard pointer 只是保持住实际需要的数据，
同时仍允许更新继续进行。

如果该读者无法合理地转换为使用引用计数，
\cref{sec:together:Long-Duration Accesses Two}
中的技巧可能会有所帮助。

% @@@ papers to maybe cite: OrcGC, ThreadScan, Fast and Robust Memory...

% @@@ Generalized hazard-pointer link counts, if and when.

% @@@ Representative hazard pointer for list, so that nothing
% @@@ in list gets freed until list's hazard pointer is released.
% @@@ Midpoint between hazard pointers and RCU, in fact, you
% @@@ could argue that Tasks Trace RCU with read-side memory
% @@@ barriers is sort of a per-CPU hazard pointers implementing RCU.
% @@@ Except no re-checking because CPUs cannot be freed.

% together/seqlock.tex
% mainfile: ../perfbook.tex
% SPDX-License-Identifier: CC-BY-SA-3.0

\section{顺序锁（sequence lock）特技}
\label{sec:together:Sequence-Locking Specials}
%
\epigraph{不会跳舞的姑娘说乐队演奏得不好。}
	 {意第绪谚语}

本节探讨 sequence lock 的一些特殊用法。

\subsection{对决的顺序锁}
\label{sec:together:Dueling Sequence Locks}

经典的 sequence lock 用例使读者能够看到
一小组变量的一致快照，
例如时间保持的校准常数。
这在实践中运行得很好，
因为校准常数很少更新，
而且更新时速度很快。
因此读者几乎不需要重试。

然而，如果更新者在更新期间被延迟，
读者也将被延迟。
这种延迟可能由中断、NMI，
甚至虚拟 CPU 抢占引起。

防止更新者延迟导致读者延迟的一种方法是
维护两组校准常数。
每组依次更新，但更新频率足够高，
使读者可以充分利用任一组。
每组都有自己的 sequence lock（\co{seqlock_t} 结构）。

更新者交替更新两组，
这样被延迟的更新者最多只会延迟其中一组的读者。

每个读者首先尝试访问第一组，
但在需要重试时尝试访问第二组。
如果第二组也迫使重试，
读者从第一组重新开始。
如果更新者卡住了，两组中只有一组会迫使读者重试，
因此读者在尝试访问另一组时就会成功。

\QuickQuiz{
	为什么不是所有 sequence lock 用例
	都以这种方式复制数据？
}\QuickQuizAnswer{
	如果数据太大，这种复制就不切实际，
	例如 \cref{sec:together:Correlated Data Elements}
	中描述的 Schr\"odinger 动物园示例。

	如果延迟被阻止，则不需要这种复制，
	例如当更新者在裸金属硬件上运行时禁用中断
	（即不使用容易发生 vCPU 抢占的虚拟机管理程序）。

	或者，如果读者可以容忍偶尔的延迟，
	则同样不需要复制。
	考虑读写锁的例子，
	写者总是延迟读者，反之亦然。

	然而，如果要复制的数据相当小，
	如果延迟是可能的，
	并且读者无法容忍这些延迟，
	复制数据是一种出色的方法。
}\QuickQuizEnd

\subsection{关联数据元素}
\label{sec:together:Correlated Data Elements}

假设我们有一个 hash table，
我们需要对其中两个或多个元素的关联视图。
这些元素一起更新，
我们不希望看到第一个元素的旧版本
与其他元素的新版本混在一起。
例如，Schr\"odinger 决定将其大家庭连同所有动物
添加到其内存数据库中。
虽然 Schr\"odinger 理解婚姻和离婚不会瞬间发生，
但他也是一个传统主义者。
因此，他绝对不希望其数据库显示
新娘已婚但新郎未婚的情况，反之亦然。
另外，如果你觉得 Schr\"odinger 是传统主义者，
试试和他的一些家庭成员交谈！
换句话说，Schr\"odinger 希望能够对其数据库
进行婚姻一致性遍历。

一种方法是使用 sequence lock
（参见 \cref{sec:defer:Sequence Locks}），
使与婚姻相关的更新在
\co{write_seqlock()} 保护下进行，
而需要婚姻一致性的读取在
\co{read_seqbegin()} / \co{read_seqretry()} 循环内进行。
请注意，sequence lock 并非 RCU 保护的替代品：
sequence lock 防止并发修改，
但仍需要 RCU 来防止并发删除。

当关联元素数量少、读取这些元素的时间短
且更新速率低时，这种方法效果很好。
否则，更新可能发生得太快，读者可能永远无法完成。
虽然 Schr\"odinger 不认为即使他最不理智的亲戚
也会快速到足以使婚姻和离婚成为问题，
但他确实意识到这个问题在其他情况下可能会出现。
避免读者饥饿（starvation）问题的一种方法是让读者在重试次数过多时
使用更新侧原语，
但这会降低性能和可扩展性。
避免\IX{starvation}的另一种方法是使用多个 sequence lock，
在 Schr\"odinger 的情况下，也许每个物种一个。

此外，如果更新侧原语使用过于频繁，
由于锁竞争（lock contention）将导致性能和可扩展性下降。
避免此问题的一种方法是维护每元素 sequence lock，
并在更新婚姻状态时持有双方配偶的锁。
读者可以在任一配偶的锁上进行重试循环，
以获得涉及该对成员的婚姻状态变化的稳定视图。
这避免了由于高结婚和离婚率造成的竞争，
但增加了在单次数据库扫描中
获取所有婚姻状态稳定视图的复杂性。

如果元素分组是明确定义且持久的
（婚姻状态被希望如此），
那么一种方法是向数据元素添加指针，
将给定组的成员链接在一起。
读者然后可以遍历这些指针来访问
与第一个找到的元素在同一组中的所有数据元素。

这种技术在 Linux 内核中被大量使用，
最值得注意的也许是 dcache 子系统~\cite{NeilBrown2015RCUwalk}。
请注意，类似的方案也可能适用于 hazard pointer。

这种方法为成功的读者提供顺序一致性，
每个读者要么看到给定更新的效果，要么看不到，
任何部分更新都会导致读侧重试。
顺序一致性是一种极强的保证，
带来同样强的限制和同样高的开销。
在这种情况下，我们看到读者可能一方面被饥饿，
另一方面可能需要获取更新侧锁。
虽然这在更新不频繁的情况下效果很好，
但它不必要地强制读侧重试，
即使更新不影响重试读者访问的任何数据。
\Cref{sec:together:Correlated Fields} 因此介绍了一种
弱得多的一致性形式，
不仅避免了读者饥饿，还避免了任何形式的读侧重试。
下一节则介绍一种更弱的一致性形式，
可以以更低的读者饥饿概率提供。

\subsection{原子移动}
\label{sec:together:Atomic Move}

假设单个数据元素从一个数据结构移动到另一个，
读者仅查找单个数据结构。
然而，当数据元素移动时，读者决不能同时在两个结构中看到它，
也决不能同时在两个结构中都看不到它。
同时，任何在新位置看到该元素的读者
决不能随后在旧位置看到它。
此外，移动可以通过将旧数据元素的新副本
插入目标位置来实现。

例如，考虑一个支持对读者原子重命名操作的 hash table。
扩展 Schr\"odinger 的动物园，
假设动物的名字发生变化，
例如，Schr\"odinger 传统家庭中的每位新娘
可能将姓氏改为与新郎相同。

但更改名字可能会改变哈希值，
也可能要求新娘的元素从一个哈希链移动到另一个。
上述一致性要求如果读者成功查找新名字，
则该读者随后对旧名字的任何查找必须失败。
类似地，如果读者对旧名字的查找导致查找失败，
则该读者随后对新名字的任何查找必须成功。
简而言之，给定读者不应看到新娘瞬间消失，
也不应该在新名字下查找新娘后又在旧名字下查找到她。

可以使用 \cref{sec:together:Correlated Data Elements} 中描述的
单个全局 sequence lock 来强制执行此一致性保证，
但即使对于不查找当前正在进行名字更改的新娘的读者，
这也可能导致读者饥饿。
此保证也可以通过要求读者获取每哈希链锁来强制执行，
但回顾
\cref{fig:datastruct:Read-Only Hash-Table Performance For Schroedinger's Zoo}
显示，即使对于单插槽系统，
这也会导致较差的性能和可扩展性。

另一种对读者更友好的实现方式是使用 RCU，
并在每个元素上放置一个 sequence lock。
查找给定元素的读者在其对该元素的
整个访问过程中充当 sequence lock 读者。
请注意，这些 sequence lock 操作将对每个读者的查找进行排序。

重命名元素可以大致按以下步骤进行：

\begin{enumerate}
\item	获取保护重命名操作的全局锁。
\item	分配并初始化一个具有新名字的元素副本。
\item	写获取旧名字元素上的 sequence lock，
	这有一个副作用：将此获取与后续的插入操作排序。
	对旧名字的并发查找现在将反复重试。
\item	插入具有新名字的元素副本。
	对新名字的查找现在将成功。
\item	执行 \co{smp_wmb()} 以将前面的插入与后续的删除排序。
\item	删除旧名字的元素。
	对旧名字的并发查找现在将失败。
\item	如有必要（例如锁依赖检查器要求），
	写释放 sequence lock。
\item	释放全局锁。
\end{enumerate}

因此，查找旧名字的读者将不断重试，
直到新名字可用，此时他们的最后一次重试将失败。
对新名字的任何后续查找将成功。
任何成功查找新名字的读者都能保证
对旧名字的任何后续查找将失败，可能需要经过一系列重试。

\QuickQuizSeries{%
\QuickQuizB{
	是否可以在新元素插入之前写获取其 sequence lock，
	而不是在旧元素被删除之前获取旧元素的 sequence lock？
}\QuickQuizAnswerB{
	可以，细节留给读者作为练习。

	术语 \emph{tombstone}（墓碑）有时用于指代
	sequence lock 被获取后具有旧名字的元素。
	类似地，术语 \emph{birthstone}（诞生石）有时用于指代
	sequence lock 仍被持有的具有新名字的元素。
}\QuickQuizEndB

\QuickQuizE{
	是否可以避免使用全局锁？
}\QuickQuizAnswerE{
	可以，一种方法是使用每哈希链锁。
	更新者可以获取与旧元素和新元素对应的锁，
	按地址顺序获取它们。
	在这种情况下，插入和删除操作当然需要
	避免获取和释放这些相同的每哈希链锁。
	如果重命名操作频繁，
	这种复杂性是值得的，
	当然也允许重命名操作并发执行。
}\QuickQuizEndE
}% End of \QuickQuizSeries

当然，也可以使用简单的标志
更高效地实现此过程。
然而，这可以被视为 sequence lock 的简化变体，
它依赖于给定元素的 sequence lock
永远不会被写获取超过一次这一事实。

% @@@ Reference TBD flag-for-deletion section.

\subsection{升级为写者}
\label{sec:together:Upgrade to Writer}

如 \cref{sec:defer:Quasi Reader-Writer Lock} 中所讨论的，
RCU 允许读者升级为写者。
当读者在扫描 RCU 保护的数据结构时
注意到当前元素需要更新时，
此功能非常有用。
当你尝试用 sequence lock 实现这个技巧时会发生什么？

事实证明，这种 sequence lock 技巧实际上在
Linux 内核中被使用，
例如 \path{drivers/infiniband/hw/hfi1/sdma.c}
中的 \co{sdma_flush()} 函数。
效果是注定包含的读者需要重试。
因此，当读者检测到某种需要重试的条件时，
会使用这个技巧。

% together/applyrcu.tex
% mainfile: ../perfbook.tex
% SPDX-License-Identifier: CC-BY-SA-3.0

\section{RCU 救援}
\label{sec:together:RCU Rescues}
%
\epigraph{大疑则大悟，小疑则小悟。}
	 {中国谚语}

本节展示如何将 RCU 应用于本书前面讨论的一些示例。
在某些情况下，RCU 提供更简洁的代码；
在另一些情况下提供更好的性能和可扩展性；
还有些情况下两者兼得。

\subsection{RCU 与基于每线程变量的统计计数器}
\label{sec:together:RCU and Per-Thread-Variable-Based Statistical Counters}

\Cref{sec:count:Per-Thread-Variable-Based Implementation}
描述了一种统计计数器的实现，它提供了出色的
性能，大致等同于简单递增（如 C 语言的 \co{++}
运算符），以及线性可扩展性——但仅限于通过
\co{inc_count()} 进行递增。
不幸的是，需要通过 \co{read_count()} 读取值的线程
必须获取全局锁，因此产生了高开销且可扩展性差。
基于锁的实现代码如
\cref{lst:count:Per-Thread Statistical Counters}
（\cpageref{lst:count:Per-Thread Statistical Counters}）所示。

\QuickQuiz{
	我们到底为什么首先需要那个全局锁？
}\QuickQuizAnswer{
	给定线程的 \co{__thread} 变量在该线程退出时消失。
	因此，需要将访问其他线程 \co{__thread} 变量的任何操作
	与线程退出同步。
	如果没有这种同步，访问刚退出线程的 \co{__thread} 变量
	将导致段错误。
}\QuickQuizEnd

\subsubsection{设计}

希望使用 RCU 而非 \co{final_mutex} 来保护
\co{read_count()} 中的线程遍历，
以便从 \co{read_count()} 获得出色的
性能和可扩展性，而不仅仅是从 \co{inc_count()}。
然而，我们不想在计算的总和中牺牲任何精度。
特别是，当给定线程退出时，
我们绝对不能丢失退出线程的计数，
也不能重复计算它。
这样的错误可能导致与结果的全精度相等的不准确性，
换句话说，这样的错误将使结果完全无用。
事实上，\co{final_mutex} 的目的之一就是
确保线程不会在 \co{read_count()} 执行过程中进出。

因此，如果我们要省去 \co{final_mutex}，
就需要想出其他方法来确保一致性。
一种方法是将所有先前退出线程的总计数
和指向每线程计数器的指针数组放入单个结构中。
这样的结构一旦对 \co{read_count()} 可用，
就保持不变，确保 \co{read_count()} 看到一致的数据。

\subsubsection{实现}

\begin{fcvref}[ln:count:count_end_rcu:whole]
\cref{lst:together:RCU and Per-Thread Statistical Counters}
的 \clnrefrange{struct:b}{struct:e}
展示了 \co{countarray} 结构，它包含一个
\co{->total} 字段用于先前退出线程的计数，
以及一个 \co{counterp[]} 指针数组指向每个当前运行线程的
每线程 \co{counter}。
此结构允许 \co{read_count()} 的给定执行
看到与指示的运行线程集一致的总计。

\begin{listing}
\ebresizeverb{.71}{\input{CodeSamples/count/count_end_rcu@whole.fcv}}
\caption{RCU 与每线程统计计数器}
\label{lst:together:RCU and Per-Thread Statistical Counters}
\end{listing}

\Clnrefrange{perthread:b}{perthread:e}
包含每线程 \co{counter} 变量的定义、
引用当前 \co{countarray} 结构的全局指针 \co{countarrayp}，
以及 \co{final_mutex} 自旋锁。

\Clnrefrange{inc:b}{inc:e} 展示了 \co{inc_count()}，
它与 \cref{lst:count:Per-Thread Statistical Counters} 中的相同。
\end{fcvref}

\begin{fcvref}[ln:count:count_end_rcu:whole:read]
\Clnrefrange{b}{e} 展示了 \co{read_count()}，它发生了显著变化。
\Clnref{rrl,rru} 用 \co{rcu_read_lock()} 和
\co{rcu_read_unlock()} 替代了 \co{final_mutex} 的获取和释放。
\Clnref{deref} 使用 \co{rcu_dereference()} 将当前
\co{countarray} 结构快照到局部变量 \co{cap} 中。
正确使用 RCU 将保证此 \co{countarray} 结构
至少在 \clnref{rru} 处当前 RCU 读侧临界区结束之前
一直伴随我们。
\Clnref{init} 将 \co{sum} 初始化为 \co{cap->total}，
即先前退出线程的计数总和。
\Clnrefrange{add:b}{add:e} 累加当前运行线程对应的每线程计数器，
最后，\clnref{ret} 返回总和。
\end{fcvref}

\begin{fcvref}[ln:count:count_end_rcu:whole:init]
\co{countarrayp} 的初始值由
\clnrefrange{b}{e} 的 \co{count_init()} 提供。
此函数在第一个线程创建之前运行，
其任务是分配并清零初始结构，
然后将其赋值给 \co{countarrayp}。
\end{fcvref}

\begin{fcvref}[ln:count:count_end_rcu:whole:reg]
\Clnrefrange{b}{e} 展示了 \co{count_register_thread()} 函数，
它由每个新创建的线程调用。
\Clnref{idx} 获取当前线程的索引，
\clnref{acq} 获取 \co{final_mutex}，
\clnref{set} 安装指向此线程 \co{counter} 的指针，
\clnref{rel} 释放 \co{final_mutex}。
\end{fcvref}

\QuickQuiz{
	\begin{fcvref}[ln:count:count_end_rcu:whole:reg]
	嘿！！！
	\cref{lst:together:RCU and Per-Thread Statistical Counters}
	的 \clnref{set} 修改了一个已有 \co{countarray} 结构中的值！
	你不是说这个结构一旦对 \co{read_count()} 可用就保持不变吗？？？
	\end{fcvref}
}\QuickQuizAnswer{
	我确实是这样说的。
	而且让 \co{count_register_thread()} 分配一个新结构
	是可能的，就像 \co{count_unregister_thread()} 当前所做的那样。

	但这是不必要的。
	回想一下基于内存快照的 \co{read_count()} 误差界限推导。
	因为新线程以初始 \co{counter} 值零开始，
	即使我们在 \co{read_count()} 执行过程中
	添加新线程，推导仍然成立。
	因此，有趣的是，在添加新线程时，
	此实现获得了分配新结构的效果，
	但实际上不必进行分配。
}\QuickQuizEnd

\begin{fcvref}[ln:count:count_end_rcu:whole:unreg]
\Clnrefrange{b}{e} 展示了 \co{count_unregister_thread()}，
它由每个线程在退出前调用。
\Clnrefrange{alloc:b}{alloc:e} 分配一个新的 \co{countarray} 结构，
\clnref{acq} 获取 \co{final_mutex}，\clnref{rel} 释放它。
\Clnref{copy} 将当前 \co{countarray} 的内容复制到新分配的版本中，
\clnref{add} 将退出线程的 \co{counter} 添加到新结构的 \co{->total}，
\clnref{null} 将退出线程的 \co{counterp[]} 数组元素置为 \co{NULL}。
\Clnref{retain} 保留指向当前（即将变为旧的）\co{countarray} 结构的指针，
\clnref{assign} 使用 \co{rcu_assign_pointer()} 安装新版本的 \co{countarray} 结构。
\Clnref{sync} 等待一个宽限期过去，
以便任何可能正在 \co{read_count()} 中并发执行、
从而可能持有对旧 \co{countarray} 结构引用的线程
能够退出其 RCU 读侧临界区，从而释放此类引用。
\Clnref{free} 然后可以安全地释放旧的 \co{countarray} 结构。
\end{fcvref}

\QuickQuiz{
	鉴于固定大小的 \co{counterp} 数组，
	这段代码究竟如何避免对线程数量的固定上限？？？
}\QuickQuizAnswer{
	你说得完全正确，该数组确实重新施加了固定上限。
	可以通过使用链表来跟踪线程来避免此限制，
	就像用户空间 RCU 中所做的那样~\cite{MathieuDesnoyers2012URCU}。
	为此代码做类似的事情留给读者作为练习。
}\QuickQuizEnd

\subsubsection{讨论}

\EQuickQuiz{
	哇！
	\Cref{lst:together:RCU and Per-Thread Statistical Counters}
	包含 70 行代码，而
	\cref{lst:count:Per-Thread Statistical Counters}
	只有 42 行。
	这种额外的复杂性真的值得吗？
}\EQuickQuizAnswer{
	这当然需要根据具体情况来决定。
	如果你需要一个线性可扩展的 \co{read_count()} 实现，
	那么 \cref{lst:count:Per-Thread Statistical Counters}
	中展示的基于锁的实现根本不适合你。
	另一方面，如果对 \co{read_count()} 的调用足够少，
	那么基于锁的版本更简单，因此可能更好，
	尽管大部分大小差异来自结构定义、
	内存分配和 \co{NULL} 返回检查。

	当然，更好的问题是``为什么语言不实现
	对 \co{__thread} 变量的跨线程访问？''
	毕竟，这样的实现将使锁和 RCU 的使用都不必要。
	这将反过来实现一个比
	\cref{lst:count:Per-Thread Statistical Counters}
	中展示的更简单的实现，
	同时具有
	\cref{lst:together:RCU and Per-Thread Statistical Counters}
	中展示的实现的所有可扩展性和性能优势！
}\EQuickQuizEnd

使用 RCU 使退出的线程能够等待，
直到其他线程保证不再使用退出线程的 \co{__thread} 变量。
这允许 \co{read_count()} 函数省去锁定，
从而为 \co{inc_count()} 和 \co{read_count()} 函数
提供出色的性能和可扩展性。
然而，这种性能和可扩展性的代价是代码复杂性的增加。
希望编译器和库的编写者使用用户级
RCU~\cite{MathieuDesnoyers2009URCU}
来提供对 \co{__thread} 变量的安全跨线程访问，
大大降低 \co{__thread} 变量用户所看到的复杂性。

\subsection{RCU 与可移除 I/O 设备的计数器}
\label{sec:together:RCU and Counters for Removable I/O Devices}

\Cref{sec:count:Applying Exact Limit Counters}
展示了一对用于处理可移除设备 I/O 访问计数的
虚构代码片段。
这些代码片段由于需要获取读写锁而在快速路径
（启动 I/O）上产生了高开销。

本节展示如何使用 RCU 来避免此开销。

执行 I/O 的代码与原始版本非常相似，
用 RCU 读侧临界区替代了原始版本中的
读写锁读侧临界区：

\begin{VerbatimN}[tabsize=8]
rcu_read_lock();
if (removing) {
	rcu_read_unlock();
	cancel_io();
} else {
	add_count(1);
	rcu_read_unlock();
	do_io();
	sub_count(1);
}
\end{VerbatimN}
\vspace{5pt}

RCU 读侧原语具有最小开销，因此如期加快了快速路径。

移除设备的更新代码片段如下：

\begin{fcvlabel}[ln:together:applyrcu:Removing Device]
\begin{VerbatimN}[tabsize=8,commandchars=\\\[\]]
spin_lock(&mylock);
removing = 1;
sub_count(mybias);
spin_unlock(&mylock);
synchronize_rcu();
while (read_count() != 0) {	\lnlbl[nextofsync]
	poll(NULL, 0, 1);
}
remove_device();
\end{VerbatimN}
\end{fcvlabel}

\begin{fcvref}[ln:together:applyrcu:Removing Device]
这里我们用独占自旋锁替代读写锁，
并添加 \co{synchronize_rcu()} 来等待所有 RCU 读侧临界区完成。
由于 \co{synchronize_rcu()}，
一旦到达 \clnref{nextofsync}，
我们就知道所有剩余的 I/O 已被计入。

当然，\co{synchronize_rcu()} 的开销可能很大，
但鉴于设备移除相当罕见，这通常是一个好的权衡。
\end{fcvref}

\FloatBarrier
\subsection{数组和长度}
\label{sec:together:Array and Length}

假设我们有一个 RCU 保护的可变长度数组，
如 \cref{lst:together:RCU-Protected Variable-Length Array} 所示。
数组 \co{->a[]} 的长度可以动态变化，
在任何给定时刻，其长度由字段 \co{->length} 给出。
当然，这引入了以下\IX{race condition}：

\begin{listing}
\begin{VerbatimL}[tabsize=8]
struct foo {
	int length;
	char *a;
};
\end{VerbatimL}
\caption{RCU 保护的可变长度数组}
\label{lst:together:RCU-Protected Variable-Length Array}
\end{listing}

\begin{enumerate}
\item	数组最初为 16 个字符长，因此 \co{->length} 等于 16。
\item	CPU~0 加载 \co{->length} 的值，获得值 16。
\item	CPU~1 将数组缩小到长度 8，
	并将指向新的 8 字符内存块的指针赋给 \co{->a[]}。
\item	CPU~0 从 \co{->a[]} 获取新指针，
	并在第 12 个元素中存储新值。
	因为数组只有 8 个字符，这导致
	SEGV 或（更糟的是）内存损坏。
\end{enumerate}

我们如何防止这种情况？

一种方法是小心使用内存屏障，
这在 \cref{chp:Advanced Synchronization: Memory Ordering} 中介绍。
这可以工作，但会产生读侧开销，
而且也许更糟的是需要使用显式内存屏障。

\begin{listing}
\begin{VerbatimL}[tabsize=8]
struct foo_a {
	int length;
	char a[0];
};

struct foo {
	struct foo_a *fa;
};
\end{VerbatimL}
\caption{改进的 RCU 保护的可变长度数组}
\label{lst:together:Improved RCU-Protected Variable-Length Array}
\end{listing}

更好的方法是将值和数组放入同一结构中，
如 \cref{lst:together:Improved RCU-Protected Variable-Length Array}~\cite{Arcangeli03}
所示。
分配新数组（\co{foo_a} 结构）就自动为数组长度提供了新的位置。
这意味着如果任何 CPU 获取了 \co{->fa} 的引用，
\co{->length} 保证与 \co{->a[]} 匹配。

\begin{enumerate}
\item	数组最初为 16 个字符长，因此 \co{->length} 等于 16。
\item	CPU~0 加载 \co{->fa} 的值，
	获得指向包含值 16 和 16 字节数组的结构的指针。
\item	CPU~0 加载 \co{->fa->length} 的值，获得值 16。
\item	CPU~1 将数组缩小到长度 8，
	并将指向包含 8 字符内存块的新 \co{foo_a} 结构的指针
	赋给 \co{->fa}。
\item	CPU~0 从 \co{->a[]} 获取新指针，
	并在第 12 个元素中存储新值。
	但因为 CPU~0 仍在引用包含 16 字节数组的
	旧 \co{foo_a} 结构，一切正常。
\end{enumerate}

当然，在两种情况下，CPU~1 都必须在释放旧数组之前
等待一个宽限期。

下一节将介绍此方法的更通用版本。

\subsection{关联字段}
\label{sec:together:Correlated Fields}
\OriginallyPublished{sec:together:Correlated Fields}{Correlated Fields}{Oregon Graduate Institute}{PaulEdwardMcKenneyPhD}

假设 Sch\"odinger 的每只动物用
\cref{lst:together:Uncorrelated Measurement Fields}
中所示的数据元素表示。
\co{meas_1}、\co{meas_2} 和 \co{meas_3} 字段是
定期更新的一组关联测量值。
读者必须看到来自单次测量更新的这三个值，
这一点至关重要：
如果读者看到 \co{meas_1} 的旧值但 \co{meas_2} 和
\co{meas_3} 的新值，该读者将陷入致命的混乱。
我们如何保证读者将看到这三个值的协调集合？\footnote{
	这种情况类似于
	\cref{sec:together:Correlated Data Elements}
	中描述的情况，
	不同之处在于这里读者只需要看到给定单个数据元素的
	一致视图，而不是前面
	\lcnamecref{sec:together:Correlated Data Elements}
	中所需的一组数据元素的一致视图。}

\begin{listing}
\begin{VerbatimL}[tabsize=8]
struct animal {
	char name[40];
	double age;
	double meas_1;
	double meas_2;
	double meas_3;
	char photo[0]; /* large bitmap. */
};
\end{VerbatimL}
\caption{非关联的测量字段}
\label{lst:together:Uncorrelated Measurement Fields}
\end{listing}

一种方法是分配一个新的 \co{animal} 结构，
将旧结构复制到新结构中，
更新新结构的 \co{meas_1}、\co{meas_2} 和 \co{meas_3} 字段，
然后通过更新指针用新结构替换旧结构。
这确实保证所有读者看到协调的测量值集合，
但由于 \co{->photo[]} 字段，需要复制大型结构。
这种复制可能产生不可接受的高开销。

\begin{listing}
\begin{VerbatimL}[tabsize=8]
struct measurement {
	double meas_1;
	double meas_2;
	double meas_3;
};

struct animal {
	char name[40];
	double age;
	struct measurement *mp;
	char photo[0]; /* large bitmap. */
};
\end{VerbatimL}
\caption{关联的测量字段}
\label{lst:together:Correlated Measurement Fields}
\end{listing}

另一种方法是施加一层间接引用，
如 \cref{lst:together:Correlated Measurement Fields}~\cite[Section 5.3.4]{PaulEdwardMcKenneyPhD}
所示。
当进行新的测量时，分配一个新的 \co{measurement} 结构，
填入测量值，并使用 \co{rcu_assign_pointer()}
更新 \co{animal} 结构的 \co{->mp} 字段指向这个新的
\co{measurement} 结构。
在宽限期过去后，旧的 \co{measurement} 结构可以被释放。

\QuickQuiz{
	但 \cref{lst:together:Correlated Measurement Fields}
	中展示的方法不会导致额外的缓存未命中，
	从而产生额外的读侧开销吗？
}\QuickQuizAnswer{
	确实可能。

\begin{listing}
\begin{VerbatimL}[tabsize=8]
struct measurement {
	double meas_1;
	double meas_2;
	double meas_3;
};

struct animal {
	char name[40];
	double age;
	struct measurement *mp;
        struct measurement meas;
	char photo[0]; /* large bitmap. */
};
\end{VerbatimL}
\caption{本地化的关联测量字段}
\label{lst:together:Localized Correlated Measurement Fields}
\end{listing}

	避免此缓存未命中开销的一种方法如
	\cref{lst:together:Localized Correlated Measurement Fields} 所示：
	只需在 \co{animal} 结构中嵌入一个名为 \co{meas} 的
	\co{measurement} 结构实例，
	并将 \co{->mp} 字段指向此 \co{->meas} 字段。

	然后可以按以下步骤进行测量更新：

	\begin{enumerate}
	\item	分配一个新的 \co{measurement} 结构，
		并将新的测量值放入其中。
	\item	使用 \co{rcu_assign_pointer()} 将 \co{->mp}
		指向此新结构。
	\item	等待宽限期过去，例如使用
		\co{synchronize_rcu()} 或 \co{call_rcu()}。
	\item	将新 \co{measurement} 结构中的测量值
		复制到嵌入的 \co{->meas} 字段中。
	\item	使用 \co{rcu_assign_pointer()} 将 \co{->mp}
		指回旧的嵌入 \co{->meas} 字段。
	\item	在另一个宽限期过去后，
		释放新的 \co{measurement} 结构。
	\end{enumerate}

	这种方法使用较重的更新过程来消除
	常见情况下的额外缓存未命中。
	额外的缓存未命中仅在更新实际进行中时才会发生。
}\QuickQuizEnd

这种方法使读者能够看到选定字段的关联值，
同时产生最小的读侧开销。
在读者仅查看单个数据元素的常见情况下，
这种每数据元素的一致性就已足够。

% Flag for deletion (if not already covered in the defer chapter).
% @@@ Issaquah Challenge.
% @@@ RLU & MV-RLU (Eventually the corresponding patents.)

\subsection{更新友好型遍历}
\label{sec:together:Update-Friendly Traversal}

假设需要对 hash table 中所有元素进行统计扫描。
例如，Schr\"odinger 可能希望计算其所有动物的
平均体长体重比。\footnote{
	为什么这样的量会有用？
	我也不知道！
	但分组统计通常是有用的。}
进一步假设 Schr\"odinger 愿意忽略
由于在统计扫描进行期间动物被添加到
hash table 或从中删除而导致的轻微误差。
Schr\"odinger 应该怎样控制并发？

一种方法是将统计扫描封装在 RCU 读侧临界区中。
这允许更新在不过度阻碍扫描的情况下并发进行。
特别是，扫描不会阻塞更新，反之亦然，
这使得即使面对非常高的更新速率，
也能优雅地支持包含大量元素的 hash table 扫描。

\QuickQuiz{
	但在调整 hash table 大小时，此扫描如何工作？
	在那种情况下，旧的和新的 hash table
	都不能保证包含 hash table 中的所有元素！
}\QuickQuizAnswer{
	确实，如
	\cref{sec:datastruct:Non-Partitionable Data Structures}
	中描述的可调整大小的 hash table
	在调整大小期间不能被完全扫描。
	一个简单的解决方法是在扫描时获取
	\co{hashtab} 结构的 \co{->ht_lock}，
	但这会阻止多个扫描同时进行。

	另一种方法是让更新在调整大小进行中时
	同时修改旧的和新的 hash table。
	这将允许扫描在旧 hash table 中找到所有元素。
	实现这一点留给读者作为练习。
}\QuickQuizEnd

\subsection{可扩展引用计数之二}
\label{sec:together:Scalable Reference Count Two}

假设\IX{reference count}正在成为性能或可扩展性的瓶颈。
你能做什么？

一种方法是为每个引用计数使用每 CPU 计数器，
与 \cref{chp:Counting} 中的算法类似，
特别是 \cref{sec:count:Exact Limit Counters}
中描述的精确限制计数器。
需要在这些计数器的每 CPU 和全局模式之间切换，
这导致要么递增和递减开销高
（\cref{sec:count:Atomic Limit Counter Implementation}），
要么使用 POSIX 信号
（\cref{sec:count:Signal-Theft Limit Counter Design}）。

另一种方法是使用 RCU 来调解每 CPU 和全局计数模式之间的切换。
每个更新在 RCU 读侧临界区内进行，
每个更新检查一个标志以确定是更新
每 CPU 计数器还是全局计数器。
要切换模式，更新标志，等待宽限期，
然后将任何剩余的计数从每 CPU 计数器
移到全局计数器，反之亦然。

Linux 内核在其 \co{percpu_ref} 引用计数器风格中使用了
这种 RCU 中介的方法。
使用此引用计数器的代码必须使用 \co{percpu_ref_init()}
初始化 \co{percpu_ref} 结构，
该函数接受的参数包括：指向结构的指针、
引用计数达到零时调用的函数的指针、
一组模式标志和一组 \co{kmalloc()} \co{GFP_} 标志。
正常初始化后，结构具有一个引用并处于每 CPU 模式。

模式标志通常为零，但如果计数器要从
慢速非每 CPU（即原子）模式开始，
可以包含 \co{PERCPU_REF_INIT_ATOMIC} 位。
还有一个 \co{PERCPU_REF_ALLOW_REINIT} 位，
允许给定的 \co{percpu_ref} 计数器通过调用
\co{percpu_ref_reinit()} 来重用，
而无需释放和重新分配。
无论 \co{percpu_ref} 结构如何初始化，
都可以使用 \co{percpu_ref_get()} 获取引用，
使用 \co{percpu_ref_put()} 释放引用。

在每 CPU 模式下，\co{percpu_ref} 结构
无法确定其值是否已达到零。
当需要这样的判断时，可以调用 \co{percpu_ref_kill()}。
此函数将结构切换到原子模式，
并删除由 \co{percpu_ref_init()} 调用安装的初始引用。
当然，在原子模式下，调用 \co{percpu_ref_get()} 和
\co{percpu_ref_put()} 相当昂贵，
但 \co{percpu_ref_put()} 可以知道值何时达到零。

希望了解更多 \co{percpu_ref} 信息的读者
请参阅 Linux 内核文档和源代码。

\subsection{重触发宽限期}
\label{sec:together:Retriggered Grace Periods}

没有 \co{call_rcu_cancel()}，
所以一旦 \co{rcu_head} 结构传递给 \co{call_rcu()}，
就无法收回。
必须让它保持不变，直到回调被调用。
在常见情况下，这是应该的，
因为 \co{rcu_head} 结构正走在通往释放的单行道上。

然而，有些用例将 RCU 与显式的 \co{open()} 和
\co{close()} 调用结合使用。
在 \co{close()} 调用之后，读者不应该开始
对数据结构的新访问，
但很可能有读者正在完成其遍历。
这种情况可以用通常的方式处理：
在 \co{close()} 调用之后等待一个宽限期，
然后再释放数据结构。

但如果在宽限期结束前调用了 \co{open()} 怎么办？

同样，没有 \co{call_rcu_cancel()}，
所以另一种方法是设置一个由回调函数检查的标志，
回调函数可以选择不实际释放任何东西。
问题解决了！

但如果在宽限期结束前先调用了 \co{open()}
然后又调用了 \co{close()} 怎么办？

一种方法是为标志设置第二个值，
使回调重新将自身排队。

但如果不仅有 \co{open()} 然后 \co{close()}，
还有在宽限期结束前的另一个 \co{open()} 呢？

在这种情况下，回调需要设置状态来反映
最后一个 \co{open()} 仍然有效。

\begin{figure}
\centering
\resizebox{2.5in}{!}{\includegraphics{together/retriggergp}}
\caption{重触发宽限期状态机}
\label{fig:count:Retrigger-Grace-Period State Machine}
\end{figure}

沿着这条思路继续，我们得到了
\cref{fig:count:Retrigger-Grace-Period State Machine}
中所示的状态机。
初始状态是 CLOSED，运行状态是 OPEN\@。
菱形箭头表示 \co{call_rcu()} 调用，
标有 ``CB'' 的箭头表示回调调用。

通过此状态机的正常路径遍历状态 CLOSED、
OPEN、CLOSING（调用 \co{call_rcu()}），
然后在回调被调用后回到 CLOSED。
如果在宽限期完成之前调用了 \co{open()}，
状态机遍历循环 OPEN、CLOSING（调用
\co{call_rcu()}）、REOPENING，
然后在回调被调用后回到 OPEN。
如果在宽限期完成之前调用了 \co{open()} 然后 \co{close()}，
状态机遍历循环 OPEN、CLOSING（调用
\co{call_rcu()}）、REOPENING、RECLOSING，
然后在回调被调用后回到 CLOSING。

给定一个无限的 \co{close()} 和 \co{open()} 交替调用序列，
状态机将遍历 OPEN 和 CLOSING（调用 \co{call_rcu()}），
然后交替停留在 REOPENING 和 RECLOSING 状态。
一旦宽限期结束，状态机将转换到
CLOSING 或 OPEN 状态，
取决于回调是在 RECLOSING 还是 REOPENING 状态中被调用。

\begin{listing}
\ebresizeverb{.88}{\input{CodeSamples/together/retrigger-gp@whole.fcv}}
\caption{重触发宽限期（伪代码）}
\label{lst:together:Retriggering a Grace Period}
\end{listing}

此状态机的粗略伪代码如
\cref{lst:together:Retriggering a Grace Period} 所示。
\begin{fcvref}[ln:together:retrigger-gp:whole]
五个状态如 \clnrefrange{states:b}{states:e} 所示，
当前状态保存在 \clnref{status} 的 \co{rtrg_status} 中，
由 \clnref{lock} 的锁保护。

三个 CB 转换（从状态 CLOSING、REOPENING
和 RECLOSING 发出）由
\clnrefrange{close_cb:b}{close_cb:e} 的 \co{close_cb()} 函数实现。
\Clnref{cleanup} 在转换到 CLOSED 状态时
调用用户提供的 \co{close_cleanup()}
来执行任何最终的清理操作，如释放内存。
\Clnref{call_rcu1} 包含 \co{call_rcu()} 调用，
它导致稍后转换到 CLOSED 状态。

\clnrefrange{open:b}{open:e} 的 \co{open()} 函数实现
到 OPEN、CLOSING 和 REOPENING 状态的转换，
\clnref{do_open} 调用 \co{do_open()} 函数来实现
任何所需数据结构的分配和初始化。

\clnrefrange{close:b}{close:e} 的 \co{close()} 函数
实现到 CLOSING 和 RECLOSING 状态的转换，
\clnref{do_close} 调用 \co{do_close()} 函数来执行
完成此转换可能需要的任何操作，
例如导致后续的只读遍历返回错误。
\Clnref{call_rcu2} 包含 \co{call_rcu()} 调用，
它导致稍后转换到 CLOSED 状态。
\end{fcvref}

此状态机和伪代码展示了如何在那些罕见的
需要此类语义的情况下获得 \co{call_rcu_cancel()} 的效果。

\subsection{长时间访问之二}
\label{sec:together:Long-Duration Accesses Two}

假设一个读写锁的读者持有锁的时间太长，
导致更新被过度延迟。
进一步假设该读者无法合理地转换为
使用引用计数
（否则参见 \cref{sec:together:Long-Duration Accesses}）。

如果该读者可以合理地转换为使用 RCU，
这可能解决问题。
原因是 RCU 读者不会完全阻塞更新，
而只阻塞那些更新的清理部分（包括内存回收）。
因此，如果系统有充足的内存，
将读写锁转换为 RCU 可能就足够了。

然而，转换为 RCU 并不总是足够的。
例如，代码可能在单个 RCU 读侧临界区内
遍历一个极其大的链式数据结构，
这可能极大地延长 RCU 宽限期，导致系统内存耗尽。
这些情况可以用几种不同的方式处理：
\begin{enumerate*}[(1)]
\item	使用 SRCU 代替 RCU，以及
\item	获取引用以退出 RCU 读者。
\end{enumerate*}

\subsubsection{使用 SRCU}
\label{sec:together:Use SRCU}

在 Linux 内核中，RCU 是全局的。
换句话说，内核中任何地方的长时间运行的 RCU 读者
都会延迟当前的 RCU 宽限期。
如果长时间运行的 RCU 读者正在遍历一个小数据结构，
那少量的数据就延迟了所有其他数据结构的释放，
这可能导致内存耗尽。

避免此问题的一种方法是对该长时间运行的
RCU 读者的数据结构使用 SRCU，
使用其自己的 \co{srcu_struct} 结构。
由此产生的长时间运行的 SRCU 读者
将仅延迟该 \co{srcu_struct} 结构的宽限期，
而不延迟 RCU 的宽限期，
从而避免内存耗尽。
更多详情请参见 \cref{sec:defer:RCU Linux-Kernel API} 中的 SRCU API。

不幸的是，这种方法确实有一些缺点。
首先，SRCU 读者不受优先级提升的约束，
这可能导致繁忙系统上低优先级 SRCU 读者的额外延迟。
更糟的是，定义单独的 \co{srcu_struct} 结构
减少了 RCU 更新者的数量，
这反过来增加了每个更新者的宽限期开销。
这意味着为每个当前 Linux 内核 RCU 用例
都提供其自己的 \co{srcu_struct} 结构，
可能会将系统范围的宽限期开销
乘以此类结构的数量。

因此，通常更好的做法是获取某种非 RCU 引用
以暂时退出 RCU 读侧临界区，
如下一节所述。

\subsubsection{获取引用}
\label{sec:together:Acquire a Reference}

如果 RCU 读侧临界区太长，就缩短它！

在某些情况下，这可以轻松完成。
例如，扫描静态分配的哈希桶数组的所有哈希链的代码
可以同样轻松地在各自的临界区中扫描每个哈希链。

这之所以有效，是因为哈希链通常相当短，而且是设计如此。
当遍历长链式结构时，必须有某种方式
在中间停下来并在稍后恢复。

例如，在 Linux 内核 v5.16 中，
\co{khugepaged_scan_file()} 函数使用 \co{need_resched()}
检查其他任务是否需要当前 CPU，
如果是，则调用 \co{xas_pause()} 适当调整
遍历的迭代器，然后调用 \co{cond_resched_rcu()}
让出 CPU\@。
\co{cond_resched_rcu()} 函数依次调用 \co{rcu_read_unlock()}、
\co{cond_resched()} 和最后 \co{rcu_read_lock()}，
以退出 RCU 读侧临界区来让出 CPU。

当这种简单的临界区缩短不可行时，
另一种方法是获取对象的引用，
然后退出 RCU 读侧临界区。
v6.17 内核中的 \co{txopt_get()} 函数采用了这种方法，
使用 \co{refcount_inc_not_zero()} 尝试获取引用。
如果由此产生的内存竞争过大，
另一个修复方法是使用 hazard pointer 而非显式引用计数。

当然，在可行的情况下，另一种方法是切换到
对暂时退出 RCU 读侧临界区更友好的
数据结构，如 hash table。

\QuickQuiz{
	但这如何与可调整大小的 hash table 一起工作，
	例如
	\cref{sec:datastruct:Non-Partitionable Data Structures}
	中描述的那个？
}\QuickQuizAnswer{
	在这种情况下，需要更多的注意，
	因为在我们暂时退出 RCU 读侧临界区期间，
	hash table 很可能被调整了大小。
	更糟的是，调整大小操作可以预期会释放旧的哈希桶，
	使我们指向空闲列表。

	但仅仅阻止旧哈希桶被释放是不够的。
	还必须确保那些桶继续被更新。

	处理此问题的一种方法是在每组桶上
	设置一个引用计数，初始值设为 1。
	全表扫描将在扫描开始时获取一个引用
	（但仅在引用非零时），
	并在扫描结束时释放它。
	调整大小操作将填充新桶、释放引用、
	等待宽限期，然后等待引用变为零。
	一旦引用为零，调整大小操作可以让更新者
	忘记旧的哈希桶，然后释放它。

	实际实现留给有兴趣的读者，
	他们将从此任务中获得很多见解。
}\QuickQuizEnd

\subsection{无锁配置更新}
\label{sec:together:Lockless Configuration Update}

假设你有一个全局指针，
它要么是 \co{NULL}，
要么指向一个提供配置信息的单例数据结构，
类似于
\cref{sec:defer:Minimal Insertion and Deletion} 中讨论的情况。
本节增加的变化是允许并发无锁更新。

\begin{listing}
\input{CodeSamples/formal/herd/C-MP+o-xchg-rcusync-o+o-xchg-rcusync-o+rl-o-o-rul@whole.fcv}
\caption{配置更新示例}
\label{lst:formal:Configuration-Update Example}
\end{listing}

乍一看，该节中的更新是无锁的，
因为它们使用 \co{rcu_assign_pointer()}。
然而，如果我们有并发的无锁更新，
一个更新的 \co{rcu_assign_pointer()}
可能被下一个更新悄悄覆盖，导致内存泄漏。
\begin{fcvref}[ln:formal:C-MP+o-xchg-rcusync-o+o-xchg-rcusync-o+rl-o-o-rul:whole]
避免这种情况的一种方法是使用原子 \co{xchg()} 指令，
如 \cref{lst:formal:Configuration-Update Example}
的 \clnref{P0:xchg, P1:xchg} 所示。
这样，被覆盖的 \co{xchg()} 返回的指针
由下一个 \co{xchg()} 返回，
允许调用者在 RCU 宽限期后释放它
（\clnref{P0:sync, P1:sync}）。
因为 Linux 内核内存模型不支持动态分配，
此代码通过存储值零来模拟释放。
\end{fcvref}

\QuickQuizSeries{%
\QuickQuizB{
	等等！！！
	\begin{fcvref}[ln:formal:C-MP+o-xchg-rcusync-o+o-xchg-rcusync-o+rl-o-o-rul:whole]
	\cref{lst:formal:Configuration-Update Example}
	的 \clnref{P0:free, P1:free}
	不需要检查从 \clnref{P0:xchg, P1:xchg} 返回的
	\co{NULL} 指针吗？
	\end{fcvref}
}\QuickQuizAnswerB{
	在这个教科书示例中，不需要。

	\begin{fcvref}[ln:formal:C-MP+o-xchg-rcusync-o+o-xchg-rcusync-o+rl-o-o-rul:whole]
	这是因为 \clnref{init:p} 将指针 \co{p} 初始化为
	指向变量 \co{a}，所以此指针永远不会是 \co{NULL}。
	然而，这是一个很好的问题，
	因为在大多数生产级代码中，
	确实需要 \co{NULL} 检查。
	此外，本节的第一句话暗示指针实际上可能是 \co{NULL}。
	添加 NULL 检查留给读者作为练习。
	\end{fcvref}
}\QuickQuizEndB
%
\QuickQuizE{
	我们在 \cref{lst:formal:Configuration-Update Example}
	中真的需要 \co{xchg()} 吗？
	为什么不用简单的加载后跟简单的存储？
}\QuickQuizAnswerE{
	是的，我们确实需要那个 \co{xchg()}。
	如果我们使用简单的加载后跟简单的存储，
	另一个进程可能在加载和存储之间更新指针，
	导致该进程的元素被覆盖和泄漏。

	但不要只听\emph{我}的话。
	用 \co{herd} 试试看！！！
	(\path{C-MP+o-o-rcusync-o+o-o-rcusync-o+rl-o-o-rul.litmus})
}\QuickQuizEndE
}		% End of \QuickQuizSeries

\begin{fcvref}[ln:formal:C-MP+o-xchg-rcusync-o+o-xchg-rcusync-o+rl-o-o-rul:whole]
读者然后可以作为直接的 RCU 读者编写
（\clnrefrange{P2:lock}{P2:unlock}），
在 \clnref{P2:load} 加载当前指针，
在 \clnref{P2:deref} 解引用它。

运行此示例
(\path{CodeSamples/formal/herd/C-MP+o-xchg-rcusync-o+o-xchg-rcusync-o+rl-o-o-rul.litmus})
通过 \co{herd}
（参见 \cref{sec:formal:Axiomatic Approaches}）
验证了此算法避免了释放后使用错误（\clnref{ex:1}），
并且它模拟了释放对象 \co{a}、\co{b} 和 \co{c}
中除一个以外的所有对象（\clnref{ex:2}），
从而也避免了内存泄漏。
\end{fcvref}

\subsection{无锁双重检查锁定}
\label{sec:together:Lockless Double-Checked Locking}

双重检查锁定是一种众所周知的模式，
提供共享状态的首次使用初始化。
但考虑一个极端低延迟环境，
其中禁止任何和所有锁的使用。
在这种约束下如何实现首次使用初始化（和后续更新）？

\begin{listing}
\ebresizeverb{.9}{
\input{CodeSamples/formal/herd/C-double-check-rcu-2@whole.fcv}
}
\caption{两进程双重检查 RCU}
\label{lst:formal:Two-Process Double-Check RCU}
\end{listing}

一种方法是使用 RCU 来保护查找操作，
并使用原子 \co{xchg()} 操作来执行更新。
与上一节一样，被覆盖的对象通过 RCU 释放
以避免干扰并发读者，例如
\cref{lst:formal:Two-Process Double-Check RCU}
(\path{C-double-check-rcu-2.litmus}) 所示。\footnote{
	对两进程示例不信服的人可以查看
	三进程示例
	(\path{C-double-check-rcu-3.litmus})。
	对三进程示例不信服的人可以构建更大的示例，
	但请注意在 Paul 的笔记本电脑上，
	\co{herd} 分析四进程示例大约需要 20 分钟，
	而三进程示例大约需要 2 秒。}

快速路径覆盖了初始化已经完成的常见情况。
\begin{fcvref}[ln:formal:C-double-check-rcu-2:whole:P0]
\Clnref{lock} 开始 RCU 读侧临界区，
\clnref{load} 获取当前指针。
如果 \clnref{if} 确定此指针非 \co{NULL}，
则初始化已完成，
\clnref{deref} 获取元素的内容，
\clnref{unl:1} 退出临界区。

否则，\clnref{if} 将看到初始化前的 \co{NULL} 指针，
并将控制转移到 \clnref{unl:2}，退出 RCU 读侧临界区。
\Clnref{init:a} 初始化元素（抽象掉任何所需的分配），
\clnref{xchg} 将指向此元素的指针与全局指针 \co{p}
原子交换，将旧值返回到 \co{r1} 中。
如果 \clnref{if:2} 确定 \co{r1} 非 \co{NULL}，
则 \clnref{sync} 等待任何访问旧元素的读者完成，
\clnref{free} 模拟 \co{kfree()}。
无论如何，\co{r2} 被设置为初始值 1。
\end{fcvref}

\begin{fcvref}[ln:formal:C-double-check-rcu-2:whole]
\Clnrefrange{P1:b}{P1:e} 以相同方式操作，
但使用元素 \co{b} 而非 \co{a}。

\Clnref{loc} 显示所有变量以便于调试 litmus 测试。
\Clnref{ex:1} 验证两个进程都获得了初始化值，
\clnref{ex:2} 验证没有内存泄漏且至少有一个元素仍在使用中。
运行 \co{herd} 确认此代码满足这些约束。
\end{fcvref}

\QuickQuiz{
	但在这样的低延迟环境中，
	内存分配不也是被禁止的吗？
}\QuickQuizAnswer{
	确实可能。
	在这种情况下，可能需要预分配，
	例如在 Linux 内核中，
	每个 CPU 可能被提供一个预分配的元素。
}\QuickQuizEnd

\subsection{无锁双重检查初始化}
\label{sec:together:Lockless Double-Checked Initialization}

假设我们有
\cref{sec:together:Lockless Double-Checked Locking}
中描述的情况，但不需要后续更新。
在这种约束下实现首次使用初始化（但不需要后续更新）
的最佳方法是什么？

\begin{listing}
\ebresizeverb{.9}{
\input{CodeSamples/formal/herd/C-double-check-cas-2@whole.fcv}
}
\caption{两进程双重检查 CAS}
\label{lst:formal:Two-Process Double-Check CAS}
\end{listing}

在这种情况下，我们只需设置指针 \co{p} 一次。
这意味着如果任何想要初始化的进程发现
此指针非 \co{NULL}，它可以简单地放弃其初始化尝试。
这反过来意味着任何读者只会获得对
唯一（永生）元素的引用，
消除了对任何延迟回收的需要。
因此，可以使用简单的 \co{cmpxchg()} 来实现此初始化过程，
如 \cref{lst:formal:Two-Process Double-Check CAS} 所示。

\begin{fcvref}[ln:formal:C-double-check-cas-2:whole:P0]
\Clnref{load} 获取指针，如果 \clnref{if} 看到它
非 \co{NULL}，\clnref{deref} 获取值。

否则，\clnref{init:a} 初始化一个元素，
\clnref{cmpxchg} 尝试将其安装在指针 \co{p} 处。
如果此尝试失败，意味着其他进程已经完成了初始化，
所以 \clnref{free} 释放此进程的元素。
\end{fcvref}

\QuickQuiz{
	\begin{fcvref}[ln:formal:C-double-check-cas-2:whole:P0]
	\cref{lst:formal:Two-Process Double-Check CAS}
	的 \clnref{load} 上的 \co{READ_ONCE()}
	不需要是 \co{smp_load_acquire()} 以防止
	与初始化并发运行的早期读者
	看到初始化前的垃圾数据吗？
	\end{fcvref}
}\QuickQuizAnswer{
	不需要，因为 \co{READ_ONCE()} 是 volatile 操作，
	因此可以作为地址依赖链的头部。
	理论上，另一个替代方案是 \co{rcu_dereference_raw()}。
	除了这可能会产生误导，
	因为该元素永远不会被释放，
	不受 \co{rcu_read_lock()} 保护，
	因此不是真正的 RCU 保护指针。

	有关地址依赖的更多信息，请参阅
	\cref{sec:memorder:Address- and Data-Dependency Difficulties}
	（\cpageref{sec:memorder:Address- and Data-Dependency Difficulties}
	开始）。
}\QuickQuizEnd

此示例强调了彻底理解需求的必要性。
在此示例中，删除运行时更新的需求
通过消除对 RCU 的需求简化了同步。\footnote{
	感谢 Denis Yaroshevskiy 指出这一点。}

\QuickQuiz{
	给定如
	\cref{lst:formal:Two-Process Double-Check CAS}
	所示的算法，
	为什么还有人使用双重检查锁定？
}\QuickQuizAnswer{
	如果标志和数据被仔细放入同一缓存行，
	那么双重检查锁定可以提供比双重检查 CAS
	更好的缓存局部性。
	这是因为在双重检查 CAS 情况下，
	指针和被指向的元素几乎总是在不同的缓存行中，
	可能导致两次缓存未命中，
	而双重检查锁定只有一次缓存未命中。

	然而，随着元素大小的增长，这个优势会减弱。
}\QuickQuizEnd

% together/microopt.tex
% mainfile: ../perfbook.tex
% SPDX-License-Identifier: CC-BY-SA-3.0

\section{微优化}
\label{sec:together:Micro-Optimization}
%
\epigraph{魔鬼藏在细节中。}{佚名}

本书中展示的数据结构都是直接编码的，
没有针对底层系统的缓存层次结构进行调整。
此外，许多实现使用函数指针来进行
键到哈希的转换和其他频繁操作。
虽然这种方法提供了简洁性和可移植性，
但在许多情况下确实会损失一些性能。

以下各节涉及特化、内存节约和硬件考虑。
请不要把这些简短的章节误认为
是对此主题的权威论述。
整本书都可以写关于针对特定 CPU 的优化，
更不用说当今常用的各种 CPU 系列了。

\subsection{特化}
\label{sec:together:Specialization}

\cref{sec:datastruct:Non-Partitionable Data Structures}
（\cpageref{sec:datastruct:Non-Partitionable Data Structures}）
中介绍的可调整大小的 hash table 使用了
不透明类型作为键。
这允许极大的灵活性，允许使用任何类型的键，
但由于通过函数指针调用也会产生显著的开销。
现代硬件使用复杂的分支预测技术来最小化此开销，
但另一方面，现实世界的软件通常比
今天大型硬件分支预测表所能容纳的更大。
对于通过指针的调用尤其如此，
在这种情况下，分支预测硬件除了
分支是否跳转的信息外，还必须记录一个指针。

可以通过将 hash table 实现特化为给定的
键类型和哈希函数来消除此开销，
例如使用 C++ 模板。
这样做可以消除
\cref{lst:datastruct:Resizable Hash-Table Data Structures}
（\cpageref{lst:datastruct:Resizable Hash-Table Data Structures}）
中 \co{ht} 结构中的 \co{->ht_cmp()}、\co{->ht_gethash()} 和
\co{->ht_getkey()} 函数指针。
它还消除了通过这些指针的相应调用，
这可以允许编译器内联生成的固定函数，
不仅消除了调用指令的开销，
还消除了参数编组的开销。

\QuickQuiz{
	这些特化实际上能节省多少？
	它们真的值得吗？
}\QuickQuizAnswer{
	第一个问题的答案留给读者作为练习。
	尝试特化可调整大小的 hash table，
	看看能带来多少性能提升。
	第二个问题无法笼统地回答，
	而必须针对特定的用例来回答。
	某些用例对性能和可扩展性极为敏感，
	而其他用例则不那么敏感。
}\QuickQuizEnd

抛开这些不谈，与我在 1970 年代初刚开始
学习编程时可用的硬件相比，
现代硬件的一大优势是所需的特化要少得多。
这比在四千字节地址空间时代
允许了更高的生产力。

\subsection{位与字节}
\label{sec:together:Bits and Bytes}

本章讨论的 hash table 几乎没有尝试节省内存。
例如，
\cref{lst:datastruct:Resizable Hash-Table Data Structures}
（\cpageref{lst:datastruct:Resizable Hash-Table Data Structures}）
中 \co{ht} 结构的 \co{->ht_idx} 字段
始终只有零或一的值，却占用了完整的 32 位内存。
可以通过从 \co{->ht_resize_key} 字段
窃取一个位来消除它。
这是可行的，因为 \co{->ht_resize_key} 字段
大到足以寻址内存中的每个字节，
而 \co{ht_bucket} 结构超过一个字节长，
所以 \co{->ht_resize_key} 字段必定有
几个多余的位可用。

这种位打包技巧经常用于高度复制的数据结构中，
例如 Linux 内核中的 \co{page} 结构。
然而，可调整大小的 hash table 的 \co{ht} 结构
并没有那么高度复制。
我们应该关注的反而是 \co{ht_bucket} 结构。
缩小 \co{ht_bucket} 结构有两个主要机会：
\begin{enumerate*}[(1)]
\item 将 \tco{->htb_lock} 字段放入
\tco{->htb_head} 指针之一的低位中，以及
\item 减少所需的指针数量。
\end{enumerate*}

第一个机会可以利用 Linux 内核中的位自旋锁，
它由 \path{include/linux/bit_spinlock.h}
头文件提供。
这些在 Linux 内核中空间关键的数据结构中使用，
但并非没有缺点：

\begin{enumerate}
\item	它们比传统的自旋锁原语慢得多。
\item	它们无法参与 Linux 内核中的
	lockdep 死锁（deadlock）检测工具~\cite{JonathanCorbet2006lockdep}。
\item	它们不记录锁的所有权，
	进一步增加了调试的复杂性。
\item	它们不参与 \rt\ 内核中的优先级提升，
	这意味着持有位自旋锁时必须禁用抢占，
	这可能降低实时延迟。
\end{enumerate}

尽管有这些缺点，
当内存紧张时，位自旋锁极其有用。

第二个机会的一个方面在
\cref{sec:datastruct:Other Resizable Hash Tables}
（\cpageref{sec:datastruct:Other Resizable Hash Tables}）
中已有介绍，该节展示了只需要一组桶列表指针
的可调整大小的 hash table，
取代了
\cref{sec:datastruct:Non-Partitionable Data Structures}
（\cpageref{sec:datastruct:Non-Partitionable Data Structures}）
中介绍的可调整大小 hash table 所需的两组。
另一种方法是使用单链桶列表
来取代本章使用的双链表。
这种方法的一个缺点是删除需要额外开销，
要么通过标记待删除的指针以便稍后删除，
要么通过搜索桶列表找到要删除的元素。

简而言之，在最小内存开销与性能和简洁性之间
存在权衡。
幸运的是，现代系统上相对较大的内存
允许我们优先考虑性能和简洁性
而非内存开销。
然而，尽管 2022 年的袖珍型智能手机配备了
数千兆字节的内存，中端服务器配备了太字节级内存，
有时仍然需要采取极端措施来减少内存开销。

\subsection{硬件考虑}
\label{sec:together:Hardware Considerations}

现代计算机通常以固定大小的块
在 CPU 和主内存之间移动数据，
块大小从 32 字节到 256 字节不等。
这些块被称为\emph{\IXpl{cache line}}，
对于高性能和可扩展性极其重要，
如 \cref{sec:cpu:Overheads} 中所讨论的。
一种久经考验的扼杀性能和可扩展性的方式是
将不兼容的变量放入同一缓存行。
例如，假设一个可调整大小的 hash table 数据元素
将 \co{ht_elem} 结构与一个频繁递增的计数器
放在同一缓存行中。
频繁的递增将导致缓存行存在于
执行递增的 CPU 上，而不存在于其他任何地方。
如果其他 CPU 尝试遍历包含该元素的哈希桶列表，
它们将遭受昂贵的缓存未命中，
降低性能和可扩展性。

\begin{listing}
\begin{VerbatimL}
struct hash_elem {
	struct ht_elem e;
	long __attribute__ ((aligned(64))) counter;
};
\end{VerbatimL}
\caption{64 字节缓存行的对齐}
\label{lst:together:Alignment for 64-Byte Cache Lines}
\end{listing}

在具有 64 字节缓存行的系统上，
解决此问题的一种方法如
\cref{lst:together:Alignment for 64-Byte Cache Lines} 所示。
这里使用 \GCC 的 \co{aligned} 属性
强制 \co{->counter} 和 \co{ht_elem} 结构
位于不同的缓存行中。
这将允许 CPU 以全速遍历哈希桶列表，
尽管有频繁的递增。

当然，这引出了一个问题：
``我们怎么知道缓存行是 64 字节的？''
在 Linux 系统上，此信息可以从
\path{/sys/devices/system/cpu/cpu*/cache/} 目录获得，
甚至可以让安装过程重建应用程序
以适应系统的硬件结构。
但是，如果你希望应用程序也能在各种环境中运行，
包括非 Linux 系统，这将更加困难。
此外，即使你满足于只在 Linux 上运行，
这种自修改安装也会带来验证方面的挑战。
例如，具有 32 字节缓存行的系统可能运行良好，
但在具有 64 字节缓存行的系统上，
由于\IX{false sharing}，性能可能会受到影响。

幸运的是，有一些经验法则在实践中效果相当好，
它们被收集在 1995 年的一篇论文中~\cite{BenjaminGamsa95a}。\footnote{
	经 Orran Krieger 许可，
	这里对其中一些规则进行了转述和扩展。}
第一组规则涉及重新排列结构以适应缓存布局：

\begin{enumerate}
\item	将只读数据放在远离频繁更新数据的位置。
	例如，将只读数据放在结构的开头，
	频繁更新的数据放在结尾。
	将很少访问的数据放在中间。
\item	如果结构具有多组字段，
	每组由独立的代码路径更新，
	则将这些组彼此分开。
	同样，在组之间放置很少访问的数据可能有帮助。
	在某些情况下，将每个这样的组
	放入由原始结构引用的单独结构中
	也可能是有意义的。
\item	在可能的情况下，将更新密集型数据与
	CPU、线程或任务关联。
	我们在 \cref{chp:Counting} 的
	计数器实现中看到了这条经验法则的
	几个非常有效的例子。
\item	更进一步，按每 CPU、每线程或每任务
	的方式划分数据，如
	\cref{chp:Data Ownership} 中所讨论的。
\end{enumerate}

已有一些关于基于追踪的自动化结构字段重排的
工作~\cite{Golovanevsky:2010:TDL:2174824.2174835}。
这项工作可能会减轻从多线程软件中获得
出色性能和可扩展性所需的
最费力的任务之一。

还有一组关于锁的经验法则：

\begin{enumerate}
\item	对于保护频繁修改数据的高竞争锁，
	采取以下方法之一：
	\begin{enumerate}
	\item	将锁放在与其保护数据不同的缓存行中。
	\item	使用适合高竞争的锁，
		例如排队锁。
	\item	重新设计以减少锁竞争。
		（这种方法最好，但并不总是简单的。）
	\end{enumerate}
\item	将无竞争的锁放入与其保护数据
	相同的缓存行中。
	这种方法意味着将锁带到当前 CPU 的缓存未命中
	同时也带来了其数据。
\item	使用\IXpl{hazard pointer}、RCU 或
	对于长时间临界区使用读写锁来保护只读数据。
\end{enumerate}

当然，这些是经验法则而非绝对规则。
需要一些实验来确定哪些最适用于给定的情况。

% batching to trade off latency for perf/scale.

\QuickQuizAnswersChp{qqztogether}
